\documentclass[11pt, a4paper]{article}

%=============== PAQUETES Y CONFIGURACIÓN ===============
\usepackage[utf8]{inputenc}
\usepackage[spanish]{babel}
\usepackage{amsmath}
\usepackage{geometry}
\usepackage{xcolor}
\usepackage{hyperref}
\usepackage{listings}

% Márgenes
\geometry{
    a4paper,
    total={170mm,257mm},
    left=20mm,
    top=20mm,
}

% Colores para el código y links
\definecolor{codegreen}{rgb}{0,0.6,0}
\definecolor{codegray}{rgb}{0.5,0.5,0.5}
\definecolor{codepurple}{rgb}{0.58,0,0.82}
\definecolor{backcolour}{rgb}{0.95,0.95,0.92}

% Configuración del paquete 'listings' para C++
\lstdefinestyle{mystyle}{
    backgroundcolor=\color{backcolour},
    commentstyle=\color{codegreen},
    keywordstyle=\color{magenta},
    numberstyle=\tiny\color{codegray},
    stringstyle=\color{codepurple},
    basicstyle=\footnotesize\ttfamily,
    breakatwhitespace=false,
    breaklines=true,
    captionpos=b,
    keepspaces=true,
    numbers=left,
    numbersep=5pt,
    showspaces=false,
    showstringspaces=false,
    showtabs=false,
    tabsize=2,
    language=C++
}
\lstset{style=mystyle}

% Configuración de hyperref
\hypersetup{
    colorlinks=true,
    linkcolor=blue,
    filecolor=magenta,
    urlcolor=cyan,
}

%=============== DOCUMENTO ===============
\title{Apunte de Programación Competitiva para la OCI}
\author{Mi Nombre} % Puedes poner tu nombre aquí
\date{\today}

\begin{document}

\maketitle
\tableofcontents
\newpage

%---------------------------------------------------
\section{Plantillas Básicas (Snippets)}
%---------------------------------------------------

\subsection{Plantilla con Casos de Prueba (Patata)}
\begin{lstlisting}
#include <bits/stdc++.h>
using namespace std;

typedef long long ll;
typedef string str;
#define vec vector

void fast_io() {
    ios_base::sync_with_stdio(false);
    cin.tie(NULL);
}

int main() {
    fast_io();
    int t;
    cin >> t;
    while (t--) {
        // Tu codigo para cada caso de prueba
    }
    return 0;
}
\end{lstlisting}

\subsection{Plantilla sin Casos de Prueba (Papa)}
\begin{lstlisting}
#include <bits/stdc++.h>
using namespace std;

typedef long long ll;
typedef string str;
#define vec vector

void fast_io() {
    ios_base::sync_with_stdio(false);
    cin.tie(NULL);
}

int main() {
    fast_io();
    // Tu codigo
    return 0;
}
\end{lstlisting}

%---------------------------------------------------
\section{Estructuras de Datos Esenciales}
%---------------------------------------------------

\subsection{\texttt{std::vector}}
Contenedor dinámico similar a un arreglo.
\begin{lstlisting}
vector<int> v;
v.push_back(10);          // Agrega al final - O(1) amortizado
v.pop_back();             // Elimina del final - O(1)
int tam = v.size();       // Obtiene el tamano
bool vacio = v.empty();   // Verifica si esta vacio
int primer = v.front();   // Primer elemento
int ultimo = v.back();    // Ultimo elemento
v.clear();                // Elimina todos los elementos

// Acceso con iteradores
auto it = v.begin();      // Iterador al inicio
auto it_end = v.end();    // Iterador al final (posicion despues del ultimo)

// Inicializacion
vector<int> v1(10, 0);    // Vector de 10 ceros
vector<int> v2 = {1,2,3}; // Inicializacion directa
\end{lstlisting}

\subsection{\texttt{std::map} y \texttt{std::unordered\_map}}
Asocian una clave a un valor.
\begin{itemize}
	\item \texttt{map}: Ordenado por clave. Operaciones en $O(\log N)$.
	\item \texttt{unordered\_map}: Sin orden. Operaciones en $O(1)$ en promedio. Más rápido.
\end{itemize}
\begin{lstlisting}
map<string, int> m;
m["hola"] = 1;
m["mundo"] = 2;

if (m.count("hola")) { // Verifica si la clave existe
    cout << m["hola"];
}
m.erase("mundo"); // Elimina por clave

for (auto const& [clave, valor] : m) {
    cout << clave << ": " << valor << endl;
}
\end{lstlisting}

\subsection{\texttt{std::set} y \texttt{std::unordered\_set}}
Almacenan elementos únicos.
\begin{itemize}
	\item \texttt{set}: Ordenado. Operaciones en $O(\log N)$.
	\item \texttt{unordered\_set}: Sin orden. Operaciones en $O(1)$ en promedio.
\end{itemize}
\begin{lstlisting}
set<int> s;
s.insert(10);
s.insert(20);
s.insert(10); // No se inserta, ya existe
s.erase(20);
if (s.count(10)) { // Verdadero
    // ...
}
\end{lstlisting}

\subsection{\texttt{std::stack} (Pila - LIFO)}
\begin{lstlisting}
stack<int> s;
s.push(10);     // Agrega
s.pop();        // Quita el de arriba
int top = s.top(); // Consulta el de arriba
bool vacio = s.empty();
\end{lstlisting}

\subsection{\texttt{std::queue} (Cola - FIFO)}
\begin{lstlisting}
queue<int> q;
q.push(10);     // Agrega al final
q.pop();        // Quita el del frente
int frente = q.front(); // Consulta el del frente
\end{lstlisting}

\subsection{\texttt{std::priority\_queue} (Cola de Prioridad)}
Por defecto, es una cola de máximos (el mayor tiene más prioridad).
\begin{lstlisting}
// Cola de maximos (por defecto)
priority_queue<int> pq_max;
pq_max.push(30); pq_max.push(10); // Top es 30

// Cola de minimos (el menor tiene mas prioridad)
priority_queue<int, vector<int>, greater<int>> pq_min;
pq_min.push(30); pq_min.push(10); // Top es 10
\end{lstlisting}


%---------------------------------------------------
\section{Algoritmos Comunes de la Librería Estándar}
%---------------------------------------------------

\subsection{\texttt{std::sort}}
Ordena un contenedor. Complejidad $O(N \log N)$.
\begin{lstlisting}
vector<int> v = {3, 1, 4, 1, 5, 9};
sort(v.begin(), v.end()); // Orden ascendente
sort(v.rbegin(), v.rend()); // Orden descendente

// Con comparador personalizado (lambda)
sort(v.begin(), v.end(), [](int a, int b) {
    return a > b; // Para orden descendente
});
\end{lstlisting}

\subsection{\texttt{std::lower\_bound} y \texttt{std::upper\_bound}}
Realizan búsqueda binaria en un rango ordenado. $O(\log N)$.
\begin{itemize}
	\item \texttt{lower\_bound(val)}: Devuelve un iterador al primer elemento que es \textbf{no menor} que `val`.
	\item \texttt{upper\_bound(val)}: Devuelve un iterador al primer elemento que es \textbf{mayor} que `val`.
\end{itemize}
\begin{lstlisting}
vector<int> v = {10, 20, 30, 30, 30, 40, 50};
//               0   1   2   3   4   5   6

// Encontrar el primer 30
auto it_l = lower_bound(v.begin(), v.end(), 30);
int idx_l = it_l - v.begin(); // idx_l = 2

// Encontrar el primer elemento > 30
auto it_u = upper_bound(v.begin(), v.end(), 30);
int idx_u = it_u - v.begin(); // idx_u = 5

// Contar ocurrencias de 30
int count = idx_u - idx_l; // count = 3
\end{lstlisting}

\subsection{\texttt{std::next\_permutation}}
Genera la siguiente permutación lexicográficamente mayor de una secuencia. Para generar todas las permutaciones, la secuencia debe estar ordenada inicialmente.
\begin{lstlisting}
string s = "123";
sort(s.begin(), s.end());
do {
    cout << s << endl;
} while (next_permutation(s.begin(), s.end()));
// Imprime: 123, 132, 213, 231, 312, 321
\end{lstlisting}

%---------------------------------------------------
\section{Técnicas Clave}
%---------------------------------------------------

\subsection{Sumas de Prefijos}
Permite calcular la suma de cualquier subarreglo en $O(1)$ después de un pre-cómputo en $O(N)$.
\begin{lstlisting}
// Arreglo original: arr
// Arreglo de sumas de prefijos: prefix_sum
int n = arr.size();
vector<ll> prefix_sum(n + 1, 0);

// 1. Construir el arreglo de prefijos
for (int i = 0; i < n; i++) {
    prefix_sum[i + 1] = prefix_sum[i] + arr[i];
}

// 2. Consultar la suma del rango [L, R] (inclusive)
// Los indices L y R son 0-based
ll suma_rango = prefix_sum[R + 1] - prefix_sum[L];
\end{lstlisting}

\subsection{Búsqueda Binaria (General)}
Técnica para encontrar una respuesta en un espacio de búsqueda monotónico.
\begin{lstlisting}
ll low = 0, high = 1e9, ans = -1;
while (low <= high) {
    ll mid = low + (high - low) / 2;
    if (check(mid)) { // check(x) es una funcion que devuelve true si x es una respuesta valida
        ans = mid;
        high = mid - 1; // Para minimizar la respuesta
        // low = mid + 1; // Para maximizar la respuesta
    } else {
        low = mid + 1; // Para minimizar la respuesta
        // high = mid - 1; // Para maximizar la respuesta
    }
}
// ans contiene la respuesta
\end{lstlisting}

%---------------------------------------------------
\section{Algoritmos de Grafos}
%---------------------------------------------------

\subsection{Representación: Lista de Adyacencia}
\begin{lstlisting}
int n; // numero de nodos
vector<vector<int>> adj; // Grafo no pesado
vector<vector<pair<int, int>>> adj_w; // Grafo pesado {vecino, peso}
\end{lstlisting}

\subsection{Búsqueda en Anchura (BFS)}
Encuentra el camino más corto en grafos no pesados.
\begin{lstlisting}
vector<int> dist(n + 1, -1);
queue<int> q;

dist[start_node] = 0;
q.push(start_node);

while (!q.empty()) {
    int u = q.front();
    q.pop();
    for (int v : adj[u]) {
        if (dist[v] == -1) {
            dist[v] = dist[u] + 1;
            q.push(v);
        }
    }
}
\end{lstlisting}

\subsection{Búsqueda en Profundidad (DFS)}
Para recorrido, detección de ciclos, componentes conexas, etc.
\begin{lstlisting}
vector<bool> visited(n + 1, false);

function<void(int)> dfs = [&](int u) {
    visited[u] = true;
    // procesar nodo u
    for (int v : adj[u]) {
        if (!visited[v]) {
            dfs(v);
        }
    }
};

dfs(start_node);
\end{lstlisting}

\subsection{Algoritmo de Dijkstra}
Camino más corto en grafos con pesos no negativos.
\begin{lstlisting}
vector<ll> dist(n + 1, 1e18);
priority_queue<pair<ll, int>, vector<pair<ll, int>>, greater<pair<ll, int>>> pq;

dist[start_node] = 0;
pq.push({0, start_node}); // {distancia, nodo}

while (!pq.empty()) {
    auto [d, u] = pq.top();
    pq.pop();

    if (d > dist[u]) continue;

    for (auto& edge : adj_w[u]) {
        auto [v, weight] = edge;
        if (dist[u] + weight < dist[v]) {
            dist[v] = dist[u] + weight;
            pq.push({dist[v], v});
        }
    }
}
\end{lstlisting}

\end{document}
